%!TEX TS-program = pdflatex
\documentclass[aspectratio=169,10pt]{beamer}
\usepackage[T2A]{fontenc}
\usepackage[utf8]{inputenc}
\usepackage[english,russian]{babel}
\usepackage{cmap}
\usepackage{amsmath,amssymb,mathtools,booktabs,array,etoolbox,PTSans,graphicx}
\graphicspath{{figures/}}
\renewcommand{\familydefault}{\sfdefault}
\renewcommand{\rmdefault}{\sfdefault}
\newbool{russian}\booltrue{russian}
\usepackage{HSE-theme/beamerthemeHSE}
\hypersetup{bookmarks=false}
\setbeamertemplate{blocks}[rounded][shadow=false]
\setbeamerfont{frametitle}{size=\Large}
\setbeamerfont{framesubtitle}{size=\small}
\setbeamerfont{block title}{size=\normalsize}
\setbeamerfont{block body}{size=\small}
\setlength{\abovedisplayskip}{5pt}
\setlength{\belowdisplayskip}{5pt}
\newcommand{\R}{\mathbb{R}}
\newcommand{\norm}[1]{\left\lVert #1\right\rVert}
\newcommand{\grad}{\nabla}
\newcommand{\figwide}[1]{\begin{center}\includegraphics[height=.52\textheight,keepaspectratio]{#1}\end{center}}
\newcommand{\figshort}[1]{\begin{center}\includegraphics[height=.40\textheight,keepaspectratio]{#1}\end{center}}
\title[Методы оптимизации]{Методы оптимизации\\в машинном обучении}
\subtitle{Лекция 11. Метод тяжёлого шарика}
\author[Н. Лукьяненко]{Никита Лукьяненко}
\institute{Совместная программа СМОЛ ГУ и ФКН НИУ ВШЭ}
\date{2026/27 учебный год}
\begin{document}

% Блок 1. Зачем нужна инерция, 7 минут. Слайды 1--5.
% 1
\frame[plain]{\titlepage}

% 2
\begin{frame}{Связь с пройденным материалом}
\small
В лекции 10 мы ускоряли градиентный спуск, меняя масштаб направлений при помощи матрицы $P$:
\[
x_{k+1}=x_k-\alpha P^{-1}\grad f(x_k).
\]
Сегодня рассмотрим другой способ использовать информацию, накопленную алгоритмом. Вместе с текущим градиентом будем учитывать, \textbf{куда точка двигалась на предыдущем шаге}.
\begin{block}{Главный вопрос}
Может ли память о прошлом перемещении сократить число шагов на вытянутой квадратичной функции?
\end{block}
\end{frame}

% 3
\begin{frame}{Почему обычный градиентный спуск замедляется}
\small
Для $q(x)=\tfrac12x^\top Ax-b^\top x$ с $A\succ0$ движение вдоль собственного вектора с собственным значением $\lambda$ задаётся множителем $1-\alpha\lambda$.
Если $\mu=\lambda_{\min}(A)\ll L=\lambda_{\max}(A)$, один шаг $\alpha$ должен одновременно подходить направлениям с малой и большой кривизной.
\[
\alpha_{\rm GD}^{*}=\frac{2}{L+\mu},\qquad
\rho_{\rm GD}^{*}=\frac{L-\mu}{L+\mu}
=\frac{\kappa-1}{\kappa+1},\quad \kappa=\frac L\mu .
\]
При большом $\kappa$ множитель $\rho_{\rm GD}^{*}$ близок к единице. Ошибка в пологом направлении сокращается медленно.
\end{frame}

% 4
\begin{frame}{Правило метода тяжёлого шарика}
\small
\[
\boxed{x_{k+1}=x_k-\alpha\grad f(x_k)+\beta(x_k-x_{k-1}).}
\]
Здесь $\alpha>0$ задаёт влияние текущего градиента, а $\beta\ge0$ сохраняет часть предыдущего перемещения $x_k-x_{k-1}$.
\begin{block}{Почему «тяжёлый шарик»}
Если точка уже двигалась в некотором направлении, член $\beta(x_k-x_{k-1})$ поддерживает это движение. Градиент изменяет скорость, а не задаёт всё перемещение заново.
\end{block}
При $\beta=0$ мы получаем обычный градиентный спуск. Название \emph{метод Поляка} также относится к этому алгоритму.
\end{frame}

% 5
\begin{frame}{Маршрут лекции: 80 минут}
\small
\begin{tabular}{@{}p{.72\linewidth}r@{}}
\toprule
Тема & Время\\
\midrule
Мотивация и связь с предыдущими лекциями & 7 мин\\
Формула метода, скорость и первые примеры & 14 мин\\
Квадратичная задача: рекурсия и устойчивость & 16 мин\\
Выбор параметров и ускоренная скорость & 20 мин\\
Что можно доказать для общей гладкой функции & 11 мин\\
Практический выбор и сравнение методов & 7 мин\\
Задача, итоги и переход к лекции 12 & 5 мин\\
\bottomrule
\end{tabular}
\end{frame}

% Блок 2. Формула и механика метода, 14 минут. Слайды 6--16.
% 6
\begin{frame}{Какие величины хранит алгоритм}
\small
В момент $k$ нам нужны текущая точка $x_k$ и предыдущая точка $x_{k-1}$. Их разность
\[
v_k=x_k-x_{k-1}
\]
будем называть \textbf{скоростью} алгоритма. Это перемещение за одну итерацию, а не физическая скорость с единицами измерения.
\[
v_{k+1}=\beta v_k-\alpha\grad f(x_k),\qquad
x_{k+1}=x_k+v_{k+1}.
\]
Такую запись удобно использовать и при вычислениях, и при доказательствах.
\end{frame}

% 7
\begin{frame}{Как запустить метод}
\small
Одной точки $x_0$ для двухшаговой формулы недостаточно: требуется начальная скорость. Обычно берут
\[
v_0=0,\qquad x_{-1}=x_0.
\]
Тогда первый шаг совпадает с градиентным:
\[
v_1=-\alpha\grad f(x_0),\qquad x_1=x_0-\alpha\grad f(x_0).
\]
Со второго шага начинает действовать инерция. Другое начальное $v_0$ возможно, но оно меняет первые итерации и константы в оценках.
\end{frame}

% 8
\begin{frame}{Одна задача, разные траектории}
\small
Пусть $q(x)=\tfrac12x^\top Ax$, где $A=R\operatorname{diag}(1,100)R^\top$, а $R$ --- поворот на $28^\circ$.
Старт $x_0=R(2,\,0.25)^\top$; для тяжёлого шарика $x_{-1}=x_0$.
ГД: $\alpha=2/101$; тяжёлый шарик: $\alpha=4/121$, $\beta=81/121$.
\begin{center}\includegraphics[height=.55\textheight,keepaspectratio]{trajectory_comparison}\end{center}
\vspace{-2mm}
Квадрат --- $x_0$, пустой круг --- $x_5$, звёздочка --- $x^*=0$. Инерция даёт более длинные шаги, включая пересечение оси долины.
\end{frame}

% 9
\begin{frame}{Откуда берётся слово «инерция»}
\small
Рассмотрим рекурсию скорости:
\[
v_{k+1}-v_k=-(1-\beta)v_k-\alpha\grad f(x_k).
\]
Левая часть показывает \textbf{изменение} перемещения. Справа есть два воздействия:
\begin{itemize}
\item $-(1-\beta)v_k$ ослабляет прежнее движение, если $0\le\beta<1$;
\item $-\alpha\grad f(x_k)$ направляет движение в сторону уменьшения функции.
\end{itemize}
При $\beta$ близком к единице прежняя скорость гасится медленно. Это помогает вдоль длинной долины, но может вызвать колебания поперёк неё.
\end{frame}

% 10
\begin{frame}{Что означает параметр $\beta$}
\small
Если градиент в нескольких последовательных точках направлен примерно одинаково, скорость накапливается:
\[
v_{k+1}=\beta v_k-\alpha g_k .
\]
При неизменном $g_k=g$ и $v_0=0$ после $k$ шагов
\[
v_k=-\alpha(1+\beta+\cdots+\beta^{k-1})g
     =-\alpha\frac{1-\beta^k}{1-\beta}g
\quad (0\le\beta<1).
\]
Такой расчёт объясняет эффект накопления. Для настоящей оптимизации градиент меняется, поэтому формула с постоянным $g$ служит только локальной иллюстрацией.
\end{frame}

% 11
\begin{frame}{Первые шаги на функции одной переменной}
\small
Пусть $f(x)=\tfrac12x^2$, $x_{-1}=x_0=1$, $\alpha=1$, $\beta=0.9$.
Поскольку $f'(x)=x$, формула принимает вид
\[
x_{k+1}=x_k-x_k+0.9(x_k-x_{k-1})
       =0.9(x_k-x_{k-1}).
\]
\begin{align*}
x_1&=0,& x_2&=-0.9,& x_3&=-0.81.
\end{align*}
Первый шаг попал точно в минимум $x^*=0$. Но накопленная скорость затем вывела точку из минимума.
\end{frame}

% 12
\begin{frame}{Почему попадание в минимум не останавливает движение}
\small
В точке $x_1=0$ градиент равен нулю, однако скорость $v_1=x_1-x_0=-1$ не равна нулю. Поэтому
\[
v_2=0.9v_1-f'(x_1)=-0.9,\qquad x_2=-0.9.
\]
\begin{alertblock}{Правило остановки}
Для метода с инерцией одного условия $\norm{\grad f(x_k)}\le\varepsilon$ может быть мало. Разумно также проверять $\norm{v_k}$ или изменение точки: иначе алгоритм способен пройти через почти стационарную точку.
\end{alertblock}
\end{frame}

% 13
\begin{frame}{Значение функции может возрасти}
\small
На предыдущем примере
\[
f(x_0)=\tfrac12,\qquad f(x_1)=0,\qquad f(x_2)=0.405.
\]
Следовательно, \textbf{обычная монотонность $f(x_{k+1})\le f(x_k)$ не является свойством метода тяжёлого шарика}.
Это само по себе не означает расходимости: алгоритм может колебаться вокруг минимума и при этом приближаться к нему.
\begin{block}{Что потребуется для анализа}
Нужно исследовать не один шаг функции, а рекурсию пары $(x_k,x_{k-1})$ либо другую величину, учитывающую скорость.
\end{block}
\end{frame}

% 14
\begin{frame}{Колебания на одномерной квадратичной функции}
\small
Задача: минимизировать $f(x)=\tfrac12x^2$; минимум $x^*=0$.
Берём $x_{-1}=x_0=1$, $\alpha=1$, $\beta=0.9$ и вычисляем
\[x_{k+1}=0.9(x_k-x_{k-1}),\qquad x_1=0,\quad x_2=-0.9.\]
\figwide{nonmonotone_one_dim}
\vspace{-2mm}
Слева показаны точки $x_k$, справа --- $f(x_k)$. Уже $f(x_2)=0.405>f(x_1)=0$: один шаг может увеличить значение функции.
\end{frame}

% 15
\begin{frame}{Два параметра решают разные задачи}
\small
\begin{tabular}{@{}p{.18\linewidth}p{.73\linewidth}@{}}
\toprule
Параметр & Роль\\
\midrule
$\alpha$ & Определяет, насколько сильно текущий градиент меняет скорость. Слишком большое значение может вызвать расходимость.\\[2mm]
$\beta$ & Сохраняет часть предыдущего перемещения. Слишком большое значение усиливает колебания и долго их гасит.\\
\bottomrule
\end{tabular}
\par\vspace{3mm}
Выбирать $\alpha$ и $\beta$ независимо по интуиции рискованно. На квадратичной задаче мы получим точную область устойчивости и совместный выбор этих параметров.
\end{frame}

% 16
\begin{frame}{Состояние метода: положение и скорость}
\small
Задача: $f(x)=\tfrac12x^2$, $x^*=0$, $x_{-1}=x_0=1$, $\alpha=0.4$, $\beta=0.7$.
На рисунке $e_k=x_k$ и $v_k=x_k-x_{k-1}$; $v_0=0$.
\[v_{k+1}=0.7v_k-0.4e_k,\qquad e_{k+1}=e_k+v_{k+1}.\]
\figshort{phase_portrait}
\vspace{-2mm}
Каждая стрелка --- один шаг. Стремление пары к $(0,0)$ означает, что $x_k\to x^*$ и перемещения затухают.
\end{frame}

% Блок 3. Квадратичная задача и устойчивость, 16 минут. Слайды 17--27.
% 17
\begin{frame}{Квадратичная функция: условия задачи}
\small
Далее до специальной оговорки рассматриваем
\[
q(x)=\tfrac12x^\top Ax-b^\top x,\qquad A=A^\top\succ0.
\]
У такой функции единственный минимум $x^*=A^{-1}b$. Обозначим
\[
\mu=\lambda_{\min}(A)>0,\qquad L=\lambda_{\max}(A),\qquad
\kappa=L/\mu .
\]
Для любого $x$ выполнено равенство $\grad q(x)=Ax-b=A(x-x^*)$: градиент \emph{точно} линейно зависит от расстояния до минимума. Поэтому формулу следующего слайда получим обычной подстановкой, без приближённого разложения.
\end{frame}

% 18
\begin{frame}{Рекурсия ошибки}
\small
Положим $e_k=x_k-x^*$. Вычтем $x^*$ из правила метода:
\begin{align*}
e_{k+1}
 &=e_k-\alpha A e_k+\beta(e_k-e_{k-1})\\
 &=\bigl((1+\beta)I-\alpha A\bigr)e_k-\beta e_{k-1}.
\end{align*}
Новая ошибка зависит от \textbf{двух} предыдущих. После перехода к собственным координатам получится числовая рекурсия второго порядка. Для неё мы проверим, какие геометрические последовательности $z^k$ являются решениями; именно эта подстановка приведёт к квадратному уравнению на слайде 21.
\end{frame}

% 19
\begin{frame}{Напоминание о собственных координатах}
\small
Так как $A$ симметрична, существует ортонормированный базис $u_1,\ldots,u_d$:
\[
Au_i=\lambda_i u_i,\qquad e_k=\sum_{i=1}^d y_{k,i}u_i.
\]
Подставляя это разложение в рекурсию ошибки, получаем по каждому $i$ отдельное уравнение. Разные координаты не смешиваются, потому что $Au_i$ остаётся вдоль $u_i$.
\begin{block}{Зачем это упрощение}
Достаточно понять одну скалярную рекурсию с произвольным $\lambda\in[\mu,L]$, а затем применить результат ко всем координатам.
\end{block}
\end{frame}

% 20
\begin{frame}{Одна скалярная рекурсия}
\small
Возьмём рекурсию со слайда 18 и разложение $e_k=\sum_i y_{k,i}u_i$.
Так как $Au_i=\lambda_i u_i$, имеем
\[
\sum_i y_{k+1,i}u_i
=\sum_i\bigl[(1+\beta-\alpha\lambda_i)y_{k,i}-\beta y_{k-1,i}\bigr]u_i.
\]
Приравняв коэффициенты при $u_i$, для координаты с $\lambda_i=\lambda$ получим
\[
\boxed{y_{k+1}=(1+\beta-\alpha\lambda)y_k-\beta y_{k-1}.}
\]
Значит, каждое собственное направление можно исследовать отдельно.
\end{frame}

% 21
\begin{frame}{Характеристический многочлен}
\scriptsize
Пусть $s_\lambda=1+\beta-\alpha\lambda$ и $\beta>0$. Для скалярной ошибки $y_{k+1}=s_\lambda y_k-\beta y_{k-1}$ проверим решение $y_k=z^k$:
\[
z^{k+1}=s_\lambda z^k-\beta z^{k-1}
\iff \boxed{p_\lambda(z)=z^2-s_\lambda z+\beta=0.}
\]
Если корни $z_+\ne z_-$, обе последовательности $z_+^k,z_-^k$ удовлетворяют рекурсии. Их линейная комбинация тоже решение. Любые $y_0,y_1$ задают коэффициенты из системы
\[
c_++c_-=y_0,\qquad c_+z_++c_-z_-=y_1
\quad(z_+\ne z_-).
\]
При двойном корне $z\ne0$: $p(z)=p'(z)=0$, поэтому $s_\lambda=2z$, $\beta=z^2$. Прямая подстановка показывает, что $kz^k$ тоже решение:
\[(k+1)z^{k+1}=2z\,kz^k-z^2(k-1)z^{k-1}.\]
Значит, общее решение $y_k=(c_1+c_2k)z^k$. При $\beta=0$ отдельно получаем ГД: $y_{k+1}=(1-\alpha\lambda)y_k$.
\end{frame}

% 22
\begin{frame}{Два типа движения}
\small
Дискриминант равен
\[
D_\lambda=(1+\beta-\alpha\lambda)^2-4\beta.
\]
\begin{itemize}
\item При $D_\lambda\ge0$ корни вещественные: ошибка вдоль данного направления складывается из двух геометрических последовательностей.
\item При $D_\lambda<0$ корни комплексно сопряжены. Их произведение равно $\beta$, поэтому модуль каждого корня равен $\sqrt{\beta}$. Ошибка колеблется с затухающей амплитудой.
\end{itemize}
Ни один тип движения сам по себе не гарантирует сходимость. Нужен модуль корней меньше единицы.
\end{frame}

% 23
\begin{frame}{Лемма о корнях квадратного уравнения}
\small
\begin{block}{Лемма}
Пусть $0\le\beta<1$ и $p(z)=z^2-sz+\beta$. Оба корня лежат строго внутри единичного круга тогда и только тогда, когда
\[
-(1+\beta)<s<1+\beta .
\]
\end{block}
Условие можно запомнить через значения многочлена на границе:
\[
p(1)=1-s+\beta>0,\qquad
p(-1)=1+s+\beta>0.
\]
На следующем слайде разберём, почему этих двух проверок достаточно при $\beta<1$.
\end{frame}

% 24
\begin{frame}{Доказательство леммы о корнях}
\small
\begin{block}{Доказательство}
Обозначим корни $z_1,z_2$. По формулам Виета
\[z_1+z_2=s,\quad z_1z_2=\beta,\quad
p(1)=(1-z_1)(1-z_2),\quad p(-1)=(1+z_1)(1+z_2).\]
Пусть $p(1)>0$ и $p(-1)>0$.
Если корни невещественные, $z_2=\overline z_1$ и
$|z_1|^2=z_1z_2=\beta<1$.
Если они вещественные, $z_1z_2=\beta\ge0$:
\begin{itemize}
\item при $s\ge0$ оба неотрицательны;
$(1-z_1)(1-z_2)>0$ даёт либо $z_1,z_2<1$, либо $z_1,z_2>1$; второй случай исключает $z_1z_2=\beta<1$;
\item при $s\le0$ оба неположительны;
$(1+z_1)(1+z_2)>0$ даёт либо $z_1,z_2>-1$, либо $z_1,z_2<-1$; второй случай снова исключает $z_1z_2<1$.
\end{itemize}
Обратно, при $|z_1|,|z_2|<1$ обе величины $p(\pm1)$ положительны (для сопряжённых корней это квадраты модулей). Из $p(1)>0$, $p(-1)>0$ получаем $-(1+\beta)<s<1+\beta$.
\end{block}
\end{frame}

% 25
\begin{frame}{Теорема: область устойчивости на квадратичной задаче}
\small
\begin{block}{Теорема}
Для $A=A^\top\succ0$ метод тяжёлого шарика сходится к $x^*$ из \emph{любой} начальной пары $(x_{-1},x_0)$ тогда и только тогда, когда
\[
\boxed{0\le\beta<1,\qquad 0<\alpha<\frac{2(1+\beta)}{L}.}
\]
\end{block}
На границах возможны корни $+1$ или $-1$, поэтому строгие неравенства существенны. Условие «из любой пары» сильнее, чем сходимость из одной специально выбранной точки.
\end{frame}

% 26
\begin{frame}{Доказательство области устойчивости}
\small
Для каждого $\lambda_i$ скалярная ошибка имеет характеристический многочлен
$p_{\lambda_i}(z)=z^2-s_{\lambda_i}z+\beta$,
$s_{\lambda_i}=1+\beta-\alpha\lambda_i$.
По лемме со слайда 24 обе его моды затухают в точности при
\[
-(1+\beta)<1+\beta-\alpha\lambda_i<1+\beta.
\]
Раскрывая два неравенства, получаем
\[
0<\alpha\lambda_i<2(1+\beta).
\]
Поскольку $\mu>0$, левая граница означает $\alpha>0$. Правая верна для всех $i$ тогда и только тогда, когда $\alpha L<2(1+\beta)$.
При $\beta\ge1$ произведение корней равно $\beta$, поэтому хотя бы один корень имеет модуль $\ge1$. Так получаем все условия теоремы.
\end{frame}

% 27
\begin{frame}{Как читать область устойчивости}
\small
По горизонтали отложено $\alpha L$, по вертикали --- $\beta$. Голубая область отвечает $0\le\beta<1$ и $0<\alpha L<2(1+\beta)$.
\figshort{stability_region}
\vspace{-2mm}
Граница показывает только сходимость на квадратичной задаче. Она ещё не говорит, насколько быстро метод сойдётся.
\end{frame}

% Блок 4. Параметры и скорость, 20 минут. Слайды 28--40.
% 28
\begin{frame}{Связь с градиентным спуском}
\small
Если положить $\beta=0$, многочлен распадается:
\[
p_\lambda(z)=z(z-(1-\alpha\lambda)).
\]
Значит, ненулевой корень равен знакомому множителю градиентного спуска $1-\alpha\lambda$. Теорема устойчивости превращается в $0<\alpha<2/L$.
\begin{block}{Результат из лекции 4}
Лучший постоянный шаг для всех $\lambda\in[\mu,L]$ равен $2/(L+\mu)$; худший модуль множителя равен $(\kappa-1)/(\kappa+1)$.
\end{block}
Теперь у нас два параметра, и мы можем уменьшить наибольший модуль корней.
\end{frame}

% 29
\begin{frame}{Что именно оптимизируем}
\small
В собственном направлении с кривизной $\lambda_i$ ошибка имеет вид
\[y_{k,i}=c_{+,i}z_+(\lambda_i)^k+c_{-,i}z_-(\lambda_i)^k,\]
а при двойном корне --- $(c_{1,i}+c_{2,i}k)z(\lambda_i)^k$.
Для геометрического затухания важно наибольшее из $|z_+|,|z_-|$: мода с меньшим модулем исчезает быстрее, а множитель $k$ не меняет предельную скорость, поскольку $k^{1/k}\to1$.
Когда известны только границы $\mu\le\lambda_i\le L$, защищаемся от худшей возможной кривизны:
\[
\rho(\alpha,\beta)=
\max_{\lambda\in[\mu,L]}
\max\{|z_+(\lambda;\alpha,\beta)|,\ |z_-(\lambda;\alpha,\beta)|\}.
\]
При $\rho<1$ все моды затухают; выбираем $\min_{\alpha>0,\,0\le\beta<1}\rho(\alpha,\beta)$.
\end{frame}

% 30
\begin{frame}{Параметры Поляка для квадратичной задачи}
\small
Пусть $\mu<L$ и $\kappa=L/\mu$. Введём
\[
r=\frac{\sqrt\kappa-1}{\sqrt\kappa+1}
=\frac{\sqrt L-\sqrt\mu}{\sqrt L+\sqrt\mu}.
\]
Тогда выбор
\[
\boxed{\alpha_*=\frac{4}{(\sqrt L+\sqrt\mu)^2},
\qquad \beta_*=r^2}
\]
даёт $\rho(\alpha_*,\beta_*)=r$. Далее мы выведем формулы и докажем, что меньшего худшего множителя постоянные параметры дать не могут.

Если $\mu=L$, то $r=0$, $\beta_*=0$ и шаг $\alpha_*=1/L$ решает квадратичную задачу за одну итерацию.
\end{frame}

% 31
\begin{frame}{Где находятся корни при этом выборе}
\small
На рисунке --- корни \textbf{характеристического уравнения для ошибки}
$p_\lambda(z)=z^2-(1+\beta_*-\alpha_*\lambda)z+\beta_*=0$ при разных $\lambda\in[\mu,L]$.
Это множители в решении $y_k$, а не точки $x_k$ алгоритма.
Здесь $\mu=1$, $L=100$, $\alpha_*=4/121$, $\beta_*=81/121$, $r=9/11$.
\begin{center}\includegraphics[height=.42\textheight,keepaspectratio]{root_circle}\end{center}
\vspace{-2mm}
$p_1(z)=(z-r)^2$, $p_{100}(z)=(z+r)^2$.
Между краями $|s_\lambda|<2r$, где $s_\lambda=1+\beta_*-\alpha_*\lambda$, поэтому
$z_\pm=(s_\lambda\pm i\sqrt{4r^2-s_\lambda^2})/2$ и $|z_\pm|=r$.
У нашей двумерной матрицы фактически лишь $\lambda=1,100$; дуги показывают возможные корни для промежуточных $\lambda$ при анализе всего интервала.
\end{frame}

% 32
\begin{frame}{Вывод параметров: крайние направления}
\scriptsize
Положим $s_\lambda=1+\beta-\alpha\lambda$. При фиксированном $\beta$ эта величина убывает от $s_\mu$ до $s_L$.
Наименьшее значение $\max\{|s_\mu|,|s_L|\}$ достигается, когда края уравновешены:
\[
s_\mu=-s_L
\quad\Longleftrightarrow\quad
\alpha=\frac{2(1+\beta)}{L+\mu}.
\]
Если $\beta=r^2$ и $|s_\lambda|\le2r$, дискриминант $s_\lambda^2-4r^2\le0$, так что оба корня имеют модуль $r$.
После уравновешивания краёв условие принимает вид
\[
\max_{\lambda\in[\mu,L]}|s_\lambda|
=\frac{(1+r^2)(L-\mu)}{L+\mu}\le2r.
\]
При наименьшем допустимом $r$ здесь равенство: оба края становятся активными,
$s_\mu=2r$, $s_L=-2r$. Это двойные корни $r$ и $-r$:
\[
1+r^2-\alpha\mu=2r,\qquad
1+r^2-\alpha L=-2r.
\]
Отсюда $\alpha\mu=(1-r)^2$, $\alpha L=(1+r)^2$,
$\kappa=((1+r)/(1-r))^2$.
Меньший худший модуль невозможен: строгое доказательство на слайдах 35--36.
\end{frame}

% 33
\begin{frame}{Почему промежуточные направления не хуже}
\small
При $\alpha_*,\beta_*$ коэффициент $s_\lambda=1+r^2-\alpha_*\lambda$ линейно убывает от $2r$ при $\lambda=\mu$ до $-2r$ при $\lambda=L$. Поэтому
\[
|s_\lambda|\le2r,\qquad
D_\lambda=s_\lambda^2-4r^2\le0.
\]
Внутри интервала корни комплексно сопряжены. Их произведение равно $r^2$, следовательно, модуль каждого равен $r$. На концах получаем двойные корни $\pm r$.
\begin{block}{Вывод}
Мы доказали верхнюю оценку $\rho(\alpha_*,\beta_*)=r$. Осталось показать, что любое другое значение $\rho$ не может быть меньше.
\end{block}
\end{frame}

% 34
\begin{frame}{Теорема: наилучший худший множитель}
\small
\begin{block}{Теорема для $A\succ0$ и $\mu<L$}
Среди постоянных $\alpha>0$ и $0\le\beta<1$ наименьшее возможное значение
\[
\max_{\lambda\in[\mu,L]}\max\{|z_+(\lambda)|,|z_-(\lambda)|\}
\]
равно
\[
\boxed{r=\frac{\sqrt\kappa-1}{\sqrt\kappa+1}}.
\]
Его достигают параметры $\alpha_*,\beta_*$ со слайда 30.
\end{block}
Это утверждение касается квадратичной функции и худшего направления в спектральном интервале. Оно не обещает такую же глобальную скорость для произвольной сильно выпуклой функции.
\end{frame}

% 35
\begin{frame}{Доказательство оптимальности: две крайние точки}
\small
Предположим, что для всех $\lambda\in[\mu,L]$ модули обоих корней не превосходят $q<1$. Тогда их произведение $\beta$ не больше $q^2$.

Для корней из круга радиуса $q$ верны $p_\mu(q)\ge0$ и $p_L(-q)\ge0$. Раскроем первое неравенство:
\[
q^2-q(1+\beta-\alpha\mu)+\beta\ge0
\quad\Longrightarrow\quad
\alpha\mu\ge(1-q)\left(1-\frac{\beta}{q}\right).
\]
Второе даёт
\[
q^2+q(1+\beta-\alpha L)+\beta\ge0
\quad\Longrightarrow\quad
\alpha L\le(1+q)\left(1+\frac{\beta}{q}\right).
\]
\end{frame}

% 36
\begin{frame}{Доказательство оптимальности: нижняя граница}
\small
Разделим верхнюю оценку для $\alpha L$ на нижнюю для $\alpha\mu$:
\[
\kappa=\frac L\mu
\le\frac{1+q}{1-q}\,
   \frac{1+\beta/q}{1-\beta/q}.
\]
Так как $\beta\le q^2$, имеем $\beta/q\le q$. Функция $(1+t)/(1-t)$ возрастает на $[0,1)$, поэтому
\[
\kappa\le\left(\frac{1+q}{1-q}\right)^2
\quad\Longrightarrow\quad
q\ge\frac{\sqrt\kappa-1}{\sqrt\kappa+1}=r.
\]
Параметры Поляка уже дали верхнюю оценку $r$. Значит, меньшее значение худшего множителя невозможно. \hfill$\square$
\end{frame}

% 37
\begin{frame}{Теорема: оценка ошибки с учётом двойных корней}
\small
\begin{block}{Теорема}
Пусть $A\succ0$, $\mu<L$, $x_{-1}=x_0$, а параметры равны $\alpha_*,\beta_*$. Тогда при всех $k\ge0$
\[
\boxed{\norm{x_k-x^*}_2
\le \bigl(1+(1+r)k\bigr)r^k\norm{x_0-x^*}_2.}
\]
Следовательно,
\[
q(x_k)-q(x^*)\le
\frac L2\bigl(1+(1+r)k\bigr)^2r^{2k}\norm{x_0-x^*}_2^2.
\]
\end{block}
Множитель $k$ нельзя просто опустить: на крайних собственных направлениях корень характеристического многочлена двойной.
\end{frame}

% 38
\begin{frame}{Доказательство оценки: одна координата}
\scriptsize
Для $\lambda\in[\mu,L]$ положим
$t=(1+r^2-\alpha_*\lambda)/(2r)\in[-1,1]$. Тогда
\[y_{k+1}=2rt\,y_k-r^2y_{k-1},\quad
y_{-1}=y_0,\quad y_1=(2rt-r^2)y_0.\]
Пусть $U_{-1}=0$, $U_0=1$ и $U_{j+1}=2tU_j-U_{j-1}$.
Последовательность $r^k(U_k-rU_{k-1})y_0$ имеет те же $y_0,y_1$ и при $k\ge1$ ту же рекурсию, поскольку
\[
r^{k+1}(U_{k+1}-rU_k)
=2rt\,r^k(U_k-rU_{k-1})-r^2r^{k-1}(U_{k-1}-rU_{k-2}).
\]
Значит, для всех $k$
\[
y_k=r^k\bigl(U_k(t)-rU_{k-1}(t)\bigr)y_0.
\]
При $t=\cos\theta$ имеем $U_j=\sin((j+1)\theta)/\sin\theta$, откуда $|U_j|\le j+1$ (и по пределу на концах).
Обозначим $C_k=r^k[(k+1)+rk]=[1+(1+r)k]r^k$.
Полученная оценка $|y_{k,i}|\le C_k|y_{0,i}|$ верна для \emph{каждой} собственной координаты. Так как базис ортонормирован,
\[
\norm{e_k}_2^2=\sum_i|y_{k,i}|^2
\le C_k^2\sum_i|y_{0,i}|^2=C_k^2\norm{e_0}_2^2.
\]
Именно это даёт первую оценку теоремы на слайде 37. Для второй используем $\lambda_i\le L$:
\[
q(x_k)-q(x^*)=\tfrac12\sum_i\lambda_i|y_{k,i}|^2
\le\tfrac L2\norm{e_k}_2^2\le\tfrac L2 C_k^2\norm{e_0}_2^2.
\]
\end{frame}

% 39
\begin{frame}{Сравнение ошибок на той же квадратичной функции}
\small
Задача та же, что на слайде 8:
$q(x)=\tfrac12x^\top R\operatorname{diag}(1,100)R^\top x$,
$R$ --- поворот на $28^\circ$, $x^*=0$, $x_0=R(2,0.25)^\top$, $x_{-1}=x_0$.
ГД: $\alpha=2/101$; тяжёлый шарик: $\alpha=4/121$, $\beta=81/121$.
\figwide{error_comparison}
\vspace{-2mm}
По вертикали $\norm{x_k-x^*}_2$ в логарифмической шкале. Оба графика вычислены по указанным рекурсиям из одной начальной точки.
\end{frame}

% 40
\begin{frame}{Почему здесь говорят об ускорении}
\small
Для большого $\kappa$ сравним главные члены:
\[
\rho_{\rm GD}^{*}=\frac{\kappa-1}{\kappa+1}
\approx 1-\frac{2}{\kappa},\qquad
r=\frac{\sqrt\kappa-1}{\sqrt\kappa+1}
\approx 1-\frac{2}{\sqrt\kappa}.
\]
\figshort{rate_vs_condition}
\vspace{-2mm}
Это скорость при больших $k$: $r<\rho_{\rm GD}^{*}$ при $\kappa>1$.
Но оценка для конкретного шага имеет вид $\norm{e_k}\le[1+(1+r)k]r^k\norm{e_0}$. Например, при $\kappa=100$ и $k=1$ её правая часть равна $\approx2.31\norm{e_0}$: эта \emph{верхняя оценка} ещё не доказывает уменьшения после первого шага. Позже $k r^k\to0$.
\end{frame}

% Блок 5. Общая гладкая функция, 11 минут. Слайды 41--47.
% 41
\begin{frame}{Чем заменить монотонность самой функции}
\scriptsize
Возьмём $f(x)=\tfrac12x^2$, $x_{-1}=x_0=1$, $\alpha=0.8$, $\beta=0.5$.
Помимо значения функции учитываем перемещение $v_k=x_k-x_{k-1}$:
\[E_k=f(x_k)+\frac{\beta}{2\alpha}v_k^2
=\underbrace{\tfrac12x_k^2}_{\text{значение функции}}
+\underbrace{0.3125v_k^2}_{\text{вклад скорости}}.\]
Это и есть \textbf{энергия} алгоритма; второй член велик, когда точка движется быстро.
Например, $x_1=0.2$, $v_1=-0.8$; $x_2=-0.36$, $v_2=-0.56$:
\[f(x_1)=0.02<f(x_2)=0.0648,\qquad E_1=0.22>E_2=0.1628.\]
\figshort{objective_and_energy}
\vspace{-2mm}
Функция выросла, но энергия уменьшилась: в этот момент запас скорости сократился.
\end{frame}

% 42
\begin{frame}{Теорема: убывание энергии}
\small
\begin{block}{Теорема}
Пусть $f$ ограничена снизу и её градиент $L$-липшицев на $\R^d$, где $L>0$. Если
\[
0\le\beta<1,\qquad 0<\alpha<\frac{2(1-\beta)}{L},
\]
то для $v_k=x_k-x_{k-1}$ энергия
\[
E_k=f(x_k)+\frac{\beta}{2\alpha}\norm{v_k}_2^2
\]
удовлетворяет
\[
E_{k+1}-E_k\le-c\norm{v_{k+1}}_2^2,\qquad
c=\frac{1-\beta}{\alpha}-\frac L2>0.
\]
Кроме того, $\grad f(x_k)\to0$.
\end{block}
Условие на шаг эквивалентно $c>0$. Это достаточное условие для общей гладкой функции; при $\beta>0$ оно строже квадратичного $\alpha<2(1+\beta)/L$ со слайда 25.
\end{frame}

% 43
\begin{frame}{Доказательство: применяем лемму об убывании}
\scriptsize
Из $L$-гладкости (лекция 4) при $x_{k+1}=x_k+v_{k+1}$:
\[
f(x_{k+1})-f(x_k)\le\grad f(x_k)^\top v_{k+1}
             +\frac L2\norm{v_{k+1}}_2^2.
\]
Подставим $\grad f(x_k)=(\beta v_k-v_{k+1})/\alpha$:
\[
f(x_{k+1})-f(x_k)\le
\frac{\beta}{\alpha}v_k^\top v_{k+1}
-\left(\frac1\alpha-\frac L2\right)\norm{v_{k+1}}_2^2.
\]
По $2a^\top b\le\norm a^2+\norm b^2$:
\[
f(x_{k+1})-f(x_k)
\le\frac{\beta}{2\alpha}\norm{v_k}^2
-\left(\frac{1-\beta/2}{\alpha}-\frac L2\right)\norm{v_{k+1}}^2.
\]
Теперь добавим разность скоростных членов энергии.
\end{frame}

% 44
\begin{frame}{Доказательство: суммируем изменения энергии}
\scriptsize
Прибавим разность скоростных членов к оценке со слайда 43:
\[
\begin{aligned}
E_{k+1}-E_k
&=f(x_{k+1})-f(x_k)
+\frac{\beta}{2\alpha}\bigl(\norm{v_{k+1}}^2-\norm{v_k}^2\bigr)\\
&\le-\left(\frac{1-\beta}{\alpha}-\frac L2\right)\norm{v_{k+1}}^2
=-c\norm{v_{k+1}}^2.
\end{aligned}
\]
Сложим это неравенство при $k=0,\ldots,N-1$:
\[
c\sum_{k=0}^{N-1}\norm{v_{k+1}}^2\le E_0-E_N\le E_0-\inf f.
\]
Правая часть не зависит от $N$, поэтому ряд квадратов сходится и $v_k\to0$.
По формуле метода
$\grad f(x_k)=(\beta v_k-v_{k+1})/\alpha\to0$.
Это вывод о стационарности; глобальная минимальность не следует.
\end{frame}

% 45
\begin{frame}{Теорема: локальная сходимость около минимума}
\small
\begin{block}{Локальная теорема}
Пусть $f\in C^2$ в окрестности $x^*$, $\grad f(x^*)=0$ и
$H_*=\grad^2f(x^*)\succ0$. Обозначим $L_*=\lambda_{\max}(H_*)$.
Если $0\le\beta<1$ и $0<\alpha<2(1+\beta)/L_*$, то при $x_{-1}$ и $x_0$ достаточно близких к $x^*$ итерации сходятся к $x^*$.
\end{block}
Если $\rho<1$ --- наибольший модуль корней квадратичной линеаризации, то для любого числа $q$ с $\rho<q<1$ найдётся $C_q$ такое, что
\[
\norm{(x_k-x^*,\,x_{k-1}-x^*)}_2
\le C_q q^k\norm{(x_0-x^*,\,x_{-1}-x^*)}_2.
\]
Обе стартовые точки должны лежать в этой окрестности.
\end{frame}

% 46
\begin{frame}{Почему локальная теорема верна}
\small
Рассмотрим отображение состояния
\[
T(x,y)=\bigl(x-\alpha\grad f(x)+\beta(x-y),\,x\bigr).
\]
Точка $(x^*,x^*)$ неподвижна. Производная $T$ в ней имеет вид
\[
DT(x^*,x^*)=
\begin{pmatrix}
(1+\beta)I-\alpha H_*&-\beta I\\ I&0
\end{pmatrix}.
\]
В собственных координатах $H_*$ её собственные значения --- корни многочленов со слайда 21. При условиях теоремы все они лежат внутри единичного круга.
\begin{block}{Последний шаг}
Для любого $q>\rho$ можно выбрать эквивалентную норму, в которой матрица производной сжимает с множителем меньше $q$. Непрерывность производной сохраняет это свойство в достаточно малой окрестности. Поэтому близкая пара состояний сходится.
\end{block}
\end{frame}

% 47
\begin{frame}{Почему теорема не утверждает глобальную сходимость}
\small
Возьмём седловую функцию $f(x,y)=\tfrac12(x^2-y^2)$. В направлении $y$ кривизна равна $-1$, и ошибка подчиняется
\[
y_{k+1}=(1+\beta+\alpha)y_k-\beta y_{k-1}.
\]
Для соответствующего многочлена $p(1)=1-(1+\beta+\alpha)+\beta=-\alpha<0$. Значит, один вещественный корень больше $1$: возмущение вдоль этого направления растёт.
\begin{alertblock}{Граница вывода}
Ускоренная оценка со слайдов 30--40 доказана для положительно определённых квадратичных функций. Локальная теорема требует положительно определённого гессиана в предельной точке и близкого старта. Для общей невыпуклой функции теорема об энергии гарантирует лишь исчезновение градиента.
\end{alertblock}
\end{frame}

% Блок 6. Выбор метода на практике, 7 минут. Слайды 48--52.
% 48
\begin{frame}{Параметры меняют и скорость, и форму траектории}
\small
На одной квадратичной функции показаны пять шагов при двух выборах $(\alpha,\beta)$. Квадрат --- старт, пустой круг --- конечная точка. Слабая инерция оставляет медленное продвижение; параметры Поляка дают более длинные шаги с заметными колебаниями.
\begin{center}\includegraphics[height=.56\textheight,keepaspectratio]{momentum_parameter_paths}\end{center}
\vspace{-2mm}
Формулы Поляка требуют оценок $\mu$ и $L$. Если эти оценки неточны, нужно проверять поведение метода и при необходимости уменьшать параметры.
\end{frame}

% 49
\begin{frame}{Как инерция сочетается с предобусловливанием}
\small
После лекции 10 можно заменить градиентный член на предобусловленный:
\[
x_{k+1}=x_k-\alpha P^{-1}\grad f(x_k)+\beta(x_k-x_{k-1}),
\qquad P\succ0.
\]
Если $P$ фиксирована, это метод тяжёлого шарика в подходящих новых координатах. На квадратичной функции его спектральный анализ проводится для
\[
B=P^{-1/2}AP^{-1/2}
\]
вместо $A$. Сначала $P$ уменьшает различие кривизн, затем инерция ускоряет движение в полученных координатах.
\begin{block}{Вычислительная цена}
Каждый шаг требует градиент и решение системы $Pp=-g$. Как и в лекции 10, полезность зависит от стоимости этого решения.
\end{block}
\end{frame}

% 50
\begin{frame}{Пример из машинного обучения: регуляризованный МНК}
\small
Возьмём два признака и два наблюдения:
$X=\operatorname{diag}(1,10)$, $y=(1,1)^\top$, коэффициент регуляризации $\lambda=1$.
Тогда
\[
f(w)=\tfrac12\norm{Xw-y}_2^2+\tfrac12\norm w_2^2,
\quad
\grad f(w)=X^\top(Xw-y)+w.
\]
Вычислим гессиан и точку минимума:
\[
A=X^\top X+I=\begin{pmatrix}2&0\\0&101\end{pmatrix},\quad
X^\top y=\binom1{10},\quad
w^*=A^{-1}X^\top y=\binom{1/2}{10/101}.
\]
Здесь $\mu=2$, $L=101$, $\kappa=101/2$. Для ГД
$\alpha_{\rm GD}=2/103$; ошибка вдоль первого признака умножается на
$1-2\alpha_{\rm GD}=99/103$.
Для тяжёлого шарика вычисляем
\[
\alpha_*=\frac4{(\sqrt{101}+\sqrt2)^2},\quad
\beta_*=\left(\frac{\sqrt{101}-\sqrt2}{\sqrt{101}+\sqrt2}\right)^2,
\quad r=\sqrt{\beta_*}.
\]
Шаг $w_{k+1}=w_k-\alpha_*\grad f(w_k)+\beta_*(w_k-w_{k-1})$ использует те же данные и дополнительно хранит $w_{k-1}$.
\end{frame}

% 51
\begin{frame}{Практическая запись алгоритма}
\small
\begin{enumerate}
\item Выберите $x_0$ и положите $v_0=0$.
\item На итерации $k$ вычислите $g_k=\grad f(x_k)$.
\item Обновите $v_{k+1}=\beta v_k-\alpha g_k$ и $x_{k+1}=x_k+v_{k+1}$.
\item Проверяйте и градиент, и величину перемещения. При больших колебаниях пересмотрите $\alpha$ или $\beta$.
\end{enumerate}
\begin{block}{Что измерять при сравнении}
Сравнивайте значения функции, норму градиента и затраченное время при одинаковом старте. На квадратичной задаче с известным $x^*$ полезна также норма ошибки $\norm{x_k-x^*}$.
\end{block}
\end{frame}

% 52
\begin{frame}{Три уже знакомые идеи ускорения}
\small
\begin{tabular}{@{}p{.25\linewidth}p{.62\linewidth}@{}}
\toprule
Метод & Что меняется в шаге\\
\midrule
Градиентный спуск & Использует только текущий градиент $g_k$.\\[1mm]
Предобусловленный GD & Меняет масштаб направления: $P^{-1}g_k$.\\[1mm]
Тяжёлый шарик & Добавляет память о перемещении $x_k-x_{k-1}$.\\
\bottomrule
\end{tabular}
\par\vspace{3mm}
В следующей лекции исследуем метод Нестерова. Он тоже вводит инерцию, но вычисляет градиент в вынесенной вперёд точке. Это различие влияет на доказуемые гарантии.
\end{frame}

% Блок 7. Закрепление, 5 минут. Слайды 53--56.
% 53
\begin{frame}{Задача для самостоятельного расчёта}
\small
Пусть $f(x)=\tfrac92x^2$, $x_{-1}=x_0=1$, $\alpha=1/9$ и $\beta=1/4$.
\begin{enumerate}
\item Вычислите $x_1$, $x_2$ и $x_3$.
\item Запишите характеристический многочлен для ошибки $x_k-x^*$.
\item Проверьте условие устойчивости и объясните, почему после попадания в минимум метод продолжает двигаться.
\end{enumerate}
\begin{block}{Подсказка}
Здесь $x^*=0$, $\lambda=L=\mu=9$ и $f'(x)=9x$. Для первого шага вспомните, что $x_{-1}=x_0$.
\end{block}
\end{frame}

% 54
\begin{frame}{Задача: подробное решение}
\small
Рекурсия имеет вид
$x_{k+1}=x_k-x_k+\tfrac14(x_k-x_{k-1})
=\tfrac14(x_k-x_{k-1})$.
Следовательно,
\[
x_1=0,\qquad x_2=-\tfrac14,\qquad x_3=-\tfrac1{16}.
\]
Многочлен:
\[
p(z)=z^2-\left(1+\tfrac14-\tfrac19\cdot9\right)z+\tfrac14
=z^2-\tfrac14z+\tfrac14.
\]
Его корни комплексно сопряжены и имеют модуль $\sqrt\beta=1/2<1$. Также
$\alpha L=1<2(1+\beta)=5/2$.
В точке $x_1=0$ градиент равен нулю, но $v_1=-1$; поэтому $x_2\ne0$.
\end{frame}

% 55
\begin{frame}{Что важно запомнить}
\small
\begin{itemize}
\item Инерция делает состояние двухшаговым: нужно отслеживать и точку, и скорость.
\item На квадратичной задаче каждое собственное направление подчиняется скалярному квадратному многочлену.
\item Для $A\succ0$ область устойчивости: $0\le\beta<1$ и $0<\alpha<2(1+\beta)/L$.
\item Параметры Поляка дают асимптотический множитель
$r=(\sqrt\kappa-1)/(\sqrt\kappa+1)$; точная оценка ошибки содержит также множитель порядка $k$.
\item Для общей гладкой функции при более осторожном шаге убывает энергия и градиент стремится к нулю.
\end{itemize}
\end{frame}

% 56
\begin{frame}{Следующая лекция: ускоренный метод Нестерова}
\small
Сегодня инерция добавлялась \emph{после} вычисления градиента в текущей точке:
\[
x_{k+1}=x_k-\alpha\grad f(x_k)+\beta(x_k-x_{k-1}).
\]
В лекции 12 мы сначала сделаем вынесенную вперёд точку
$y_k=x_k+\beta_k(x_k-x_{k-1})$, а затем вычислим градиент \emph{в ней}:
\[
x_{k+1}=y_k-\alpha_k\grad f(y_k).
\]
Сравним оба метода на тех же квадратичных примерах и разберём, почему ускорение Нестерова имеет гарантии для более широкого класса выпуклых задач.
\end{frame}

\end{document}
